\subsection{Hypocotyl dataset collected under a light--dark cycle}
\label{subsec:hypocotyldata}

Hypocotyl-length measurements were collected in this study from \textit{Arabidopsis thaliana} grown at 23~$^{\circ}$C under a 12-h light/12-h dark photoperiod (12L12D). The experiment included five genotypes: Col-0, \textit{hy5}, MLB (\textit{ML1} promoter::phyB-GFP/\textit{phyB-9}), EMS57, and the \textit{phyA phyB} double mutant. Lengths were measured in millimetres at 0, 12, 24, 36, 48, 60, and 72~h, with illumination during the first 0--12-h interval. The analysis included 943 measurements: 197 for Col-0, 196 for \textit{hy5}, 225 for MLB, 94 for EMS57, and 231 for \textit{phyA phyB}. Replicate counts ranged from 6 to 38 per genotype--time combination. The dataset therefore combined sparse temporal sampling with unequal replication and missing entries. Measurements under continuous red light were excluded from all partitions. Time was expressed as elapsed hours because developmental age at the initial observation could not be confirmed. Irradiance, spectral composition, independent growth batches, and some mutant alleles were not documented in the available records.

The dataset is intended for public release with the article, together with the original workbook, individual measurements linked to source cells, available experimental metadata, fixed evaluation partitions, and preprocessing scripts. It provides a complementary resource for modelling genotype-level growth from sparsely sampled replicate plants under a prescribed illumination schedule at constant temperature.

Longitudinal plant identities could not be verified from the workbook; observations were therefore modelled as replicate samples defining five genotype-level growth curves. Records were grouped by genotype and source row, and existing group assignments were retained after restricting the data to 12L12D. Each group was assigned to one partition across observation times, yielding 547 training, 171 validation, and 225 test measurements. Every partition contained all genotype--time combinations. Source-row grouping provided a conservative partition without assuming that a row identified the same plant at successive times. The protocol held out replicate groups within a single experiment; independent experimental batches were not established.

Training and primary evaluation used replicate means calculated independently within each partition. Let $\bar H_{s,g}(t)$ denote the mean of the available lengths for genotype $g$ at time $t$ in partition $s$. Each partition contained 35 targets, corresponding to five genotypes at seven times. Missing entries were excluded before averaging, and missing genotype--time groups contributed no target. No phenotype values were imputed. Individual measurements were retained for a secondary evaluation of prediction error.

Lengths were normalized by the maximum training measurement, $a=13.719$~mm. The data loss is
\begin{equation}
 \mathcal L_{\mathrm{data}}^{\mathrm{hyp}}
 =\frac{1}{5}\sum_{g=1}^{5}\left[
 \frac{1}{|\mathcal T_{\mathrm{train},g}|}
 \sum_{t\in\mathcal T_{\mathrm{train},g}}
 \left(\hat y_g(t)-\frac{\bar H_{\mathrm{train},g}(t)}{a}\right)^2
 \right]^{1/2},
 \label{eq:hypocotyl_data_loss}
\end{equation}
where $\mathcal T_{s,g}$ contains the observed times for genotype $g$ in partition $s$. Primary RMSE is the mean of the five genotype-specific trajectory RMSEs against $\bar H_{s,g}(t)$, expressed in millimetres. Relative RMSE normalizes this score by the mean of the evaluated targets, as in Eq.~\eqref{eq:rrmse}. Each observed time receives equal weight within its genotype trajectory, irrespective of replicate count. These metrics quantify prediction error for mean growth; individual-level error is reported separately.

\subsubsection{Light-conditioned latent dynamics and an ordinary logistic reference}
\label{subsec:lightmodel}

Logistic descriptions of hypocotyl elongation provide a parsimonious growth reference \citep{alimchandani2026atlas}. We adopt the standard zero-offset formulation,
\begin{equation}
 \frac{dH_g}{dt}=r_g H_g\left(1-\frac{H_g}{K_g}\right),
 \qquad
 H_g(t)=\frac{K_g}{1+(K_g/H_g(0)-1)\exp(-r_g t)}.
 \label{eq:hypocotyl_logistic}
\end{equation}
Each genotype has a positive, time-independent growth rate $r_g$ and a carrying capacity $K_g$. For the process-based baseline, $(r_g,K_g,H_g(0))$ were estimated separately for each genotype by fitting the training replicate means from eight deterministic starting points, giving 15 parameters. Initial length was estimated exclusively from training data. The fitted rate represents an effective growth coefficient over the observation interval; light- and dark-phase rates were not estimated separately.

Light signalling and circadian regulation influence diurnal hypocotyl growth \citep{nozue2007rhythmic}. PhytoODE represents the illumination phase through the prescribed binary input,
\begin{equation}
 L(t)=
 \begin{cases}
  1,&0\leq (t\bmod 24)<12,\\
  0,&12\leq (t\bmod 24)<24,
 \end{cases}
 \qquad t\ \text{in hours}.
 \label{eq:hypocotyl_light}
\end{equation}
Temperature was constant and was excluded from the model inputs. The prescribed illumination schedule, including its state at 72~h, was available independently of the phenotype measurements. Because all genotypes shared one schedule, $L(t)$ was a deterministic function of elapsed time. This experiment therefore assessed the utility of an explicit illumination feature within a known photoperiod; it did not evaluate transfer to new photoperiods or identify a causal light response.

The PhytoODE architecture comprised a four-dimensional genotype embedding, an encoder of width 8, and an eight-dimensional latent state. The encoder receives normalized time $\tau=t/72$ and the binary illumination schedule. The vector field contains one hidden layer of width 16 and receives the latent state, genotype embedding, normalized time, and instantaneous illumination state. A decoder of width 16 produces the scaled length $\hat y_g$, corresponding to $\hat H_g=a\hat y_g$ in physical units. An auxiliary genotype-dependent network outputs $r_g\in(0,0.5)$~h$^{-1}$ and a scaled capacity $\kappa_g\in(0,2)$, with $K_g=a\kappa_g$ in millimetres. Including the auxiliary network, the model contained 1,103 parameters, compared with 1,055 when the illumination channel was omitted.

Predicted length is determined by the latent dynamics. Equation~\eqref{eq:hypocotyl_logistic} regularizes the decoded growth rate, evaluated as
\begin{equation}
 \frac{d\hat y_g}{dt}
 =\frac{1}{72}\nabla_{\mathbf z}G_\psi(\mathbf z_g)^{\top}
 \mathbf F_\theta(\mathbf z_g,L(72\tau),\mathbf e_g,\tau).
 \label{eq:hypocotyl_derivative}
\end{equation}
Automatic differentiation evaluates the decoder gradient while retaining the computational graph through the residual and Runge--Kutta solver. The resulting loss provides gradients for the latent dynamics, decoder, encoder, genotype embedding, and auxiliary parameters without finite differences of measured lengths. Integration used a 3-h grid $\mathcal G$ containing every illumination transition. The illumination state was held constant throughout each Runge--Kutta interval, including its final stage, and updated at the start of the next interval. Predicted length remains continuous across each transition, although its derivative may change. The right-hand derivative was used at transition nodes in Eq.~\eqref{eq:hypocotyl_derivative}; this calculation does not require differentiating the binary illumination signal.

The physics residual was evaluated at all 23 interior grid nodes, denoted $\mathcal C$, independently of the observation mask:
\begin{align}
 \mathcal L_{\mathrm{ODE}}^{\mathrm{hyp}}
 &=\frac{1}{5|\mathcal C|}\sum_{g=1}^{5}\sum_{t\in\mathcal C}
 \left[\frac{d\hat y_g}{dt}
 -r_g\hat y_g\left(1-\frac{\hat y_g}{\kappa_g}\right)\right]^2,
 \label{eq:hypocotyl_residual}\\
 \mathcal L_K^{\mathrm{hyp}}
 &=\frac{1}{5}\sum_{g=1}^{5}
 \left|\kappa_g-\max_{t\in\mathcal G}\hat y_g(t)\right|.
 \label{eq:hypocotyl_capacity}
\end{align}
The training objective is $\mathcal L_{\mathrm{data}}^{\mathrm{hyp}}+\lambda_{\mathrm{ODE}}\mathcal L_{\mathrm{ODE}}^{\mathrm{hyp}}+\lambda_K\mathcal L_K^{\mathrm{hyp}}$. Illumination conditions the latent dynamics, while the logistic reference uses one time-independent $r_g$ and $\kappa_g$ per genotype. The capacity penalty enforces consistency with the maximum predicted length over the observation interval; mature hypocotyl length was not available as a separate observation.

\subsubsection{Baselines and validation-based selection}

The primary comparison comprised PhytoODE with illumination input, an unregularized latent ODE, Logistic-PINN, LSTM-NN, random forest, and the standard logistic ODE. Baseline inputs comprised genotype and time, excluding the illumination channel. The unregularized latent ODE used the same hidden dimensions as PhytoODE, with both physics coefficients set to zero and an independently selected learning rate; omission of illumination reduced its input dimension. Logistic-PINN comprised a coordinate MLP with a four-dimensional genotype embedding, normalized time, two hidden layers of width 32, and the same logistic derivative and capacity penalties. Its time derivative was evaluated by automatic differentiation of the coordinate network. This architecture differs from the LSTM-based Logi-PINN used in the temperature-conditioned benchmarks. LSTM-NN used time as its recurrent input, recurrent widths of 16 and 8, and a four-dimensional genotype embedding concatenated before a dense layer of width 16. Random forest used one-hot genotype encoding and normalized time.

Model selection for the genotype-and-time predictors evaluated 16 configurations each for PhytoODE and Logistic-PINN, combining learning rates $\{0.003,0.01\}$, derivative coefficients $\{0.5,5,50,500\}$, and capacity coefficients $\{0.01,0.1\}$. The unregularized latent ODE and LSTM-NN were evaluated at both learning rates. Random forest used 300 trees with minimum leaf sizes $\{1,2,4\}$. The two best seed-1 configurations for each stochastic model were evaluated with seeds 2 and 3, and selection used three-seed mean validation RMSE. These fitted baselines were retained for the illumination-input comparison with identical data partitions and targets.

For PhytoODE with illumination input, 32 configurations were evaluated with seed 1. Sixteen used the same coefficient and learning-rate grid, and sixteen were sampled by stratification in logarithmic space over learning rate, $\lambda_{\mathrm{ODE}}$, $\lambda_K$, and weight decay. The respective ranges were $[10^{-3},10^{-1.7}]$, $[10^{-2},10^3]$, $[10^{-4},10^{-0.5}]$, and $[10^{-6},10^{-3}]$; two sampled configurations used zero weight decay. All neural models were trained for 1,500 epochs with Adam, gradient clipping at 1, and cosine learning-rate decay. Checkpoint selection began at epoch 200 and used the mean of the three most recent validation evaluations, conducted every 20 epochs. The three best seed-1 configurations were evaluated with seeds 2 and 3. A prespecified feature comparison added illumination while retaining the hyperparameters selected for PhytoODE without illumination: $(\lambda_{\mathrm{ODE}},\lambda_K)=(50,0.01)$ and learning rate 0.01. This configuration was evaluated with all three seeds. The main PhytoODE configuration was selected by the lowest three-seed mean validation RMSE among these candidates. The illumination-input study comprised 32 initial training runs and eight additional runs for seed confirmation.

The fixed-hyperparameter comparison assessed the effect of adding illumination without additional hyperparameter selection, although the input modification also changed weight-matrix dimensions and initialization. Neither this comparison nor the baseline latent ODE isolates the physics-loss contribution to the illumination-conditioned model. The benchmark evaluates complete predictors with model-specific inputs and optimization budgets; the additional search was applied exclusively to PhytoODE. Reported standard deviations quantify variability across training seeds and do not estimate biological sampling uncertainty.

The data partitions were retained from preliminary analyses in which test results had been inspected. The illumination-input candidate set and stopping criterion were specified before training, and configurations were selected using training and validation data before test evaluation. No candidates were added in response to the resulting test errors. Nevertheless, the reused partition does not constitute an independent test cohort. The planned release will include the evaluation protocols, configurations, validation histories, and predictions.
